% Guide: Writing a Top-Tier Thesis on a Serverless Zero‑Knowledge Browser Password Manager
% Audience: undergraduate engineering/CS thesis aiming for “elite” research-thesis quality  
% Goal: help you *write the thesis* (structure, what to mention, what to prove), not write the thesis itself.

% ---

% 0) The thesis “contract” (what you are actually defending)

% 0.1 Your central research claim (make it explicit early)
% You are not “building a password manager”. You are defending a claim like:

% - **Claim A (security feasibility):** A browser extension using **native WebCrypto** can enforce **zero‑knowledge** confidentiality + integrity even when storage and transport are untrusted (“hostile storage”).
% - **Claim B (architectural resilience):** Removing application-server logic (“serverless” as *no data-management logic on the provider*) reduces systemic risk (no mass-exfiltration target), at the cost of client-side sync complexity.
% - **Claim C (practicality):** The design remains usable and performant (KDF cost, UI latency, conflict handling) for realistic vault sizes.

% Everything in the document should support, qualify, or limit these claims.

% 0.2 Your deliverable triad
% Your thesis should clearly separate these deliverables:
% 1) **Security specification** (threat model + crypto suite + invariants)  
% 2) **System design** (hexagonal architecture + components + data flows)  
% 3) **Experimental evaluation** (security testing + performance + usability + trade-offs)

% The software artifact is the apparatus used to validate the above.

% ---

% 1) Thesis outline (recommended chapter split)

% This split matches typical engineering templates (Intro / SoTA / Design / Implementation / Evaluation) but upgraded to research quality.

% Chapter 1 — Introduction (argument framing)
% **Purpose:** establish urgency + gap + research question + contributions.

% Include:
% - The credential-management problem: hundreds of services, reuse, stuffing risk.
% - Why SaaS password managers centralize risk (vendor breach → millions of vaults).
% - Why browsers matter (browser as OS) and why web crypto used to be controversial.
% - Your definition of **serverless** (no app server logic; storage is “dumb blob sync”).
% - **Research question(s)** and **hypothesis** (security + practicality).
% - **Scope:** Web Extension first; local-first via IndexedDB; sync optional; platform roadmap (addon → desktop → Android → iOS).
% - **Contributions (bulleted, measurable):**
%  - Hostile-storage threat model + security spec (suite-v1, anti-rollback, AAD metadata binding, pepper strategy)
%  - Hexagonal architecture with ports/adapters and strict dependency direction
%  - Local-first sync protocol with **mandatory manual conflict resolution**
%  - Benchmarks + security verification results + usability findings
% - **Non-goals / boundaries:** browser/OS compromise, malicious extension with high privilege, etc.

% **Rule:** no feature list. Everything must support the thesis claim(s).

% ---

% Chapter 2 — Background & Literature Review (critical synthesis, not glossary)
% **Purpose:** justify design decisions by reviewing what exists and what fails.

% What to cover:
% - **Browser security model:** SOP, XSS as bypass, CSP relevance.
% - **History of JS crypto vs native APIs:** pitfalls of JS polyfills, need for CSPRNG, WebCrypto turning point (`crypto.getRandomValues`, `crypto.subtle`).
% - **Password manager taxonomy** (table):
%  - Local (KeePass), Cloud SaaS, Browser extension SaaS-like, **Your serverless hostile-storage model**
%  - Compare: trust boundary, attack surface, metadata leakage, sync strategy.
% - **Related approaches and baselines:**
%  - Mopass-like approaches (position as comparison)
%  - AWS KMS-style designs (trade: server/provider trust and complexity)
% - **Cryptographic primitives overview** *only as needed to defend your choices*:
%  - AEAD concept, AES-GCM integrity, IV uniqueness requirement
%  - KDFs: PBKDF2 vs Argon2 (note WebCrypto native support vs WASM tradeoff)
%  - Signatures: why signing vault modifications matters for provenance

% **Rule:** Every subsection ends with “Implication for our design: …”.

% ---

% Chapter 3 — Problem Definition & Threat Model (your “security backbone”)
% **Purpose:** formalize what you protect, from whom, and what you don’t.

% 3.1 Assets (be explicit)
% - Vault plaintext entries (passwords, notes, TOTP seeds)
% - Vault key (DEK), device private keys, derived master key material
% - Metadata: vaultId, deviceId, timestamps, revisions, device registry
% - Sync configuration (bucket name, region, credentials)
% - Local caches (IndexedDB) and remote blobs (S3)

% 3.2 Adversaries & capabilities
% Use a structured model (STRIDE is fine) and then specialize into your hostile-storage assumptions:
% - **Untrusted storage:** attacker can read/write/rollback/delete objects in S3 and IndexedDB
% - **Offline attacker:** can attempt brute-force after exfiltration
% - **Compromised transport:** treat TLS as insufficient for confidentiality (corporate proxies etc.)
% - **Trusted boundary:** only extension volatile memory at runtime
% - Optional: discuss realistic extension threats (XSS within extension UI, malicious websites trying to interact, etc.)

% 3.3 Security goals (turn into testable invariants)
% Examples (write as “must” statements):
% - Confidentiality: storage reveals only ciphertext; no plaintext leaves trusted memory
% - Integrity: modified/forged vault snapshots are rejected
% - Anti-rollback: older valid snapshots cannot silently replace newer state
% - Key secrecy: vault key and device private keys are non-extractable and never stored unwrapped
% - Sync correctness: no silent overwrites; conflicts require explicit user choice

% 3.4 Out-of-scope (must be explicit)
% - Compromised browser binary/OS malware/physical memory extraction
% - Fully malicious extension with broad privileges (define whether considered out-of-scope or partially mitigated)
% - Side-channel resistance beyond WebCrypto guarantees

% ---

% Chapter 4 — Security Specification (Suite-v1) (write like an engineering spec)
% **Purpose:** define cryptographic suite, formats, metadata, and versioning.

% Include:

% 4.1 Suite versioning strategy
% - “suite-v1” is frozen; future “suite-v2” must remain backward compatible by envelope tagging.

% 4.2 Primitive selection table (with rationale)
% - AES-256-GCM for payload AEAD (12-byte IV uniqueness requirement; tag length)
% - AES-256-KW for wrapping the vault DEK (RFC 3394)
% - Ed25519 signatures for vault modifications (provenance / non-repudiation)
% - ECDH P-256 for device-to-device shared secrets during setup
% - PBKDF2-HMAC-SHA256 with **≥600k iterations** (justify with benchmarks)

% > Note: if any of these primitives are not uniformly supported in your target browsers via WebCrypto, document compatibility and fallbacks (or constrain supported browsers). Don’t hide this—own it.

% 4.3 Key hierarchy (diagram required)
% Show:
% - Master password → KDF → KEK (key-encryption key)
% - DEK (vault key) wrapped by KEK (AES-KW)
% - Device signing keypair (Ed25519), private key non-extractable
% - Device pairing key agreement (ECDH) for provisioning secrets

% 4.4 Pepper strategy (split-knowledge hardening)
% Specify precisely:
% - Input = SHA256(masterPassword + applicationPepper)
% - Pepper location: embedded in extension build (not in vault/config)
% - Threat benefit: attacker needs both storage exfiltration + client reverse engineering

% Also be honest:
% - Pepper is not a silver bullet; it raises cost and changes attacker model.

% 4.5 Integrity binding with AAD + authenticated metadata
% Define which metadata fields are bound in AES-GCM AAD:
% - vaultId, signerDeviceId, revision, timestamp (and any others)
% Explain the security property:
% - metadata tampering breaks tag verification

% 4.6 Anti-rollback logic
% - Track last-seen timestamp/revision locally
% - Reject/flag older snapshots even if signature verifies
% - Define the UX for rollback warnings (must not silently accept)

% 4.7 Runtime security & memory hygiene requirements
% Mandate:
% - Use `Uint8Array` not strings where possible
% - Explicit `.fill(0)` wiping
% - `extractable: false` for DEK and private keys
% - Auto-lock default 5 minutes triggers wipe flow

% **Rule:** treat this chapter as your *“ground truth”* for later evaluation.

% ---

% Chapter 5 — System Architecture (Hexagonal / Ports & Adapters)
% **Purpose:** demonstrate engineering rigor and testability.

% 5.1 Architectural style and why
% - Hexagonal architecture to isolate core logic from infrastructure and UI frameworks.
% - Dependency rule: always inward; core has no React/no AWS SDK.

% 5.2 Layer definitions (map to folders)
% - `src/core/`: types + ports (interfaces), pure logic
% - `src/adapters/`: WebCrypto adapter, S3 adapter, IndexedDB adapter, mocks
% - `src/ui/`: views/hooks/stores (React)
% - `extension/background/`: service worker coordinating sync (no React)

% 5.3 State management & DI rules (be strict)
% - React Context used only for **dependency injection** (wiring ports → adapters)
% - Global state with Zustand (why over Redux: simplicity/perf)
% - Business logic in hooks; views are presentational

% 5.4 File conventions (enforceable)
% - `.type.ts`, `.port.ts`, `.adapter.ts`, `.store.ts`, `.hook.ts`, `.view.tsx`
% Explain how this supports portability: swapping S3→GCS = new adapter only.

% 5.5 Core data flow diagrams (must include)
% - Unlock flow: master password → derive → unwrap DEK → load vault entries
% - CRUD flow: entry edit → encrypt → commit local → mark dirty for sync
% - Sync flow overview (ties into Chapter 6)

% ---

% Chapter 6 — Synchronization Protocol (local-first, user-controlled)
% **Purpose:** specify your distributed system behavior without a server arbiter.

% 6.1 Design principle: “No Auto-Merge for conflicts”
% Explain why LWW is unsafe in a password manager:
% - can overwrite a new password with old one
% - can resurrect deleted/compromised credential
% Your stance: **manual conflict resolution is mandatory**.

% 6.2 The 6-step sync algorithm (spec + pseudo)
% 1) initiation (preconditions: enabled, unlocked, network check)
% 2) download remote encrypted snapshot
% 3) double decrypt (local + remote in memory)
% 4) diff + categorize:
%   - auto-resolved (one-sided changes)
%   - conflicts (both modified or delete-vs-modify)
% 5) resolution UI:
%   - masked fields by default
%   - explicit per-item choice: Use Local / Use Remote / Skip
%   - no batch resolve
% 6) completion:
%   - merge → encrypt → write local
%   - atomic remote upload (temp then rename)

% 6.3 Edge cases (explicit section)
% - Offline handling (abort + queue retry)
% - Master password rotation (detect decrypt failure on other device; prompt for new pass)
% - Large vault performance (pagination + progress in Diff UI)

% **Rule:** this chapter must be detailed enough that another engineer could implement it.

% ---

% Chapter 7 — Multi-Device Setup, Identity, and Revocation
% **Purpose:** solve provisioning without a central server.

% 7.1 What must be shared vs what must never be transferred
% - Must share: storage config + encryption salt
% - Must never transfer electronically: master password (manual entry only)

% 7.2 Provisioning methods (3 methods, with security properties)
% 1) QR transfer (co-located):
%   - encrypted session payload, expires in 5 minutes, one-time use
% 2) Encrypted export file:
%   - config encrypted with key derived from master password
%   - safe over email/USB/cloud drive (without password memory)
% 3) Manual entry

% 7.3 Device identity model
% - unique signerDeviceId (UUID)
% - Ed25519 signing key per device
% - device registry metadata (OS/browser/last sync), opt-in to reduce fingerprinting

% 7.4 Revocation semantics
% - cannot wipe disconnected device (no server)
% - rotate vault keys to prevent future updates from revoked device

% ---

% Chapter 8 — Data Model, Organization, and UX Semantics (security-relevant UX)
% **Purpose:** show that usability is part of security, and organization decisions are research-informed.

% Include:
% - “Place vs Context” ontology:
%  - Folders: hierarchical tree, mutually exclusive, icon-based
%  - Tags: flat grouped list, multi-select, color-based
% - Global library of 100+ standard definitions (bundled JSON)
% - Archetypes (Developer / Family / Standard) to seed structure
% - Smart suggestions (normalize names/icons/colors)

% Tie to outcomes:
% - reduces over-organization
% - improves consistency and retrieval (security via reduced mistakes)

% ---

% Chapter 9 — Infrastructure Backends & Portability (S3 as reference, not dependency)
% **Purpose:** defend “serverless” operational simplicity and provider-agnostic design.

% Cover:
% - S3 reference implementation
% - IaC: CloudFormation template provisions:
%  - S3 bucket (versioning, public access blocked)
%  - Cognito Identity Pool
%  - IAM least privilege (GetObject/PutObject in prefix)

% Alternative backends (and what changes):
% - GCS (auth differences → adapter)
% - Azure Blob (mapping differences)
% - MinIO (S3-compatible endpoint)
% - Storj (requires distinct adapter)

% Include:
% - cost analysis for 1MB vault daily sync (pennies/month, often free-tier)

% ---

% Chapter 10 — Implementation Notes (evidence of engineering, not a code tour)
% **Purpose:** demonstrate you implemented the spec and faced real constraints.

% Include:
% - repo overview (structure, build, testing)
% - WebCrypto specifics:
%  - async flow patterns to avoid UI blocking
%  - IV generation and uniqueness policy
%  - key wrapping/unwrapping
% - storage choices:
%  - avoid localStorage, use IndexedDB for blobs and async transactions
% - “war stories”:
%  - JS GC + memory wiping limits
%  - how you constrained sensitive lifetimes
%  - extension MV3 service worker constraints (if relevant)

% **Rule:** every implementation subsection references the requirement/spec it satisfies.

% ---

% Chapter 11 — Evaluation (the grade is decided here)
% **Purpose:** validate claims quantitatively and adversarially.

% Split into three pillars:

% 11.1 Security verification (“negative testing”)
% - XSS simulation attempts:
%  - try exporting keys; prove `extractable: false` blocks export
% - hostile storage tampering:
%  - modify ciphertext / metadata; ensure GCM tag failure
%  - rollback attempt with older snapshot; ensure detection + warning
% - network interception:
%  - show only high-entropy blobs transmitted (no plaintext)
% - automated audits:
%  - Lighthouse / OWASP ZAP security headers (CSP, HSTS, etc.)
%  - document findings and fixes

% Deliverables:
% - test cases table: threat → test → expected result → observed result

% 11.2 Performance benchmarking
% Define test environment (browser version, CPU, RAM).

% Measure:
% - PBKDF2 latency at multiple iteration counts (include 600k)
% - encrypt/decrypt throughput (WebCrypto vs pure JS lib baseline if used)
% - diff+merge performance for 1k+ entries, UI responsiveness

% Include:
% - graphs and tables
% - justification for chosen KDF iterations based on usability/security tradeoff

% 11.3 Usability study
% n=10–20:
% - SUS score
% - task success time: sign up, save credential, autofill
% - trust perception: “browser-based vs desktop” sentiment
% - analyze feedback specifically about:
%  - manual conflict resolution burden
%  - setup complexity (S3 provisioning)

% ---

% Chapter 12 — Discussion, Limitations, Future Work
% **Purpose:** interpret results honestly, not “victory lap”.

% Include:
% - whether claims A/B/C held, and under what conditions
% - limitations:
%  - browser/OS compromise
%  - pepper strategy limits (reverse engineering)
%  - manual conflict UX cost
% - future work:
%  - WebAuthn hardware-backed key material replacing master password
%  - post-quantum experiments via WASM (Kyber), clearly marked as future work
%  - improved secure enclave usage where available (platform constraints)

% ---

% 2) What diagrams/tables you must include (minimum set)

% 1) **Threat model table** (adversary → capability → mitigation → residual risk)  
% 2) **Key hierarchy diagram** (master → KEK → wrap DEK; device signing keys; pairing keys)  
% 3) **Envelope format / metadata fields** (what is AAD; what is signed; version tags)  
% 4) **Hexagonal architecture diagram** (ports & adapters; dependency direction)  
% 5) **Sync sequence diagram** (6-step flow, diff UI pause, atomic upload)  
% 6) **Comparison table** vs SaaS PMs, Mopass, KMS-like approaches  
% 7) **Benchmark tables/plots** (PBKDF2 iterations vs latency; crypto throughput; diff performance)  
% 8) **Usability results** (SUS + task times + qualitative themes)

% Rule: every figure/table must be referenced and explained in text.

% ---

% 3) Writing rules (to match “elite thesis” expectations)

% 3.1 Narrative flow
% Each chapter must answer:
% - What problem remains?
% - What insight/design choice addresses it?
% - How you validate it?

% 3.2 Tense discipline
% - Past: what you implemented/tested (“was measured”, “was designed”)
% - Present: established facts (“AES-GCM provides…”)

% 3.3 Precision (ban “very secure”)
% Replace vague claims with measurable statements:
% - “resists tampering” → “GCM tag verification fails on any bitflip; modification is rejected”
% - “hard to brute force” → “PBKDF2 600k iterations yields X ms on test machine; attacker cost model…”

% 3.4 Citations strategy
% Cite for:
% - password reuse / credential stuffing prevalence
% - arguments about JS crypto pitfalls and WebCrypto motivation
% - local-first sync and conflict-resolution patterns (where applicable)
% - usability instruments (SUS)
% - cryptographic standards (W3C WebCrypto, RFCs)

% ---

% 4) “Do not forget” checklist (common thesis-killers)

% - Define **serverless** precisely (no server-side data logic; storage is dumb).
% - Make **hostile storage** assumption explicit (remote and local stores attacker-modifiable).
% - State **what is trusted** (only volatile memory of extension instance).
% - Specify **suite-v1** and the migration plan to suite-v2.
% - Document IV uniqueness policy (AES-GCM requirement).
% - Explain **pepper** benefits *and* limits.
% - Bind metadata with **AAD**, and add anti-rollback logic.
% - Explain **manual conflict resolution** and why auto-merge is rejected.
% - Show portability: adapters, provider-agnostic design, MinIO/Storj feasibility.
% - Evaluate with adversarial tests + benchmarks + usability (not just “it works”).

% ---

% 5) Suggested work plan (so the writing stays consistent)

% 1) Freeze **Threat Model + Suite-v1 spec** (Ch. 3–4) first  
% 2) Then write **Architecture + Sync spec** (Ch. 5–6)  
% 3) Implement to match the spec  
% 4) Write **Evaluation plan** before running tests (so it’s not cherry-picked)  
% 5) Run tests, capture artifacts (screenshots of ZAP findings, benchmark logs, diff cases)  
% 6) Write Discussion/Limitations with the same rigor as the design chapters

% ---

% 6) Thesis “quality bar” rubric you can self-grade against

% - **Originality:** serverless redefinition + hostile-storage + no-auto-merge protocol + pepper + AAD metadata binding
% - **Technical depth:** spec is implementable; architecture is testable; crypto is correctly used
% - **Critical thinking:** trade-offs are explicit (setup complexity, manual conflicts, out-of-scope threats)
% - **Evaluation rigor:** negative testing + reproducible benchmarks + usability study
% - **Professionalism:** diagrams, tables, consistent terminology, strong citations, no filler screenshots

% ---

% If you want, paste your current table of contents (or your university template headings), and I’ll map them 1:1 to the guide above (what to put in each section, what evidence/figures to include, and what mistakes to avoid).